\documentclass[10pt,a4paper]{amsart}
\usepackage[margin=19mm]{geometry}
\usepackage{amsmath,amssymb,mathtools}
\usepackage[scr=rsfs,cal=euler]{mathalfa}
\usepackage{xfrac}
\usepackage[unicode,hidelinks,bookmarks=false]{hyperref}
\newcommand{\ie}{i.e.\ }
\newcommand{\cf}{cf.\ }
\newcommand{\resp}{respectively\ }
\newcommand{\Subsection}{\subsection}
\providecommand{\nobreakdash}{\nobreak\hbox{-}}
\setlength{\emergencystretch}{2em}
\multlinegap0pt
\usepackage{bm}
\usepackage{bbm}
\usepackage{dsfont}
\usepackage{xspace}
\usepackage{cancel}
\usepackage{mathtools}
\usepackage{amsthm}
\makeatletter
\@ifundefined{dfn}{\newtheorem{dfn}{Dfn}}{}
\@ifundefined{thm}{\newtheorem{thm}[dfn]{Theorem}}{}
\makeatother
\title[Irreducible representations of tree automorphism groups into]{Irreducible representations of tree automorphism groups into Pontryagin spaces}
\author{Federico Viola}
\date{Source: arXiv:2508.02367v2 (CC BY 4.0)}
\begin{document}
\maketitle

\noindent\emph{Source and licence.} Irreducible representations of tree automorphism groups into Pontryagin spaces. Version arXiv:2508.02367v2 (CC BY 4.0), distributed under the Creative Commons Attribution 4.0 International licence, as recorded in the arXiv licence metadata for this exact version and shown on the abstract page https://arxiv.org/abs/2508.02367v2. Full rights notice: licenses/arxiv-2508.02367-CC-BY-4.0.txt.\par\smallskip

\noindent\textit{Editorial scope.} Counted unit: the Thm whose source label is \texttt{sph2} in the article (source0.tex lines 385--387, source label \texttt{sph2}), delivered together with the article's own complete original proof of it (source0.tex lines 389--417, one \texttt{proof} environment of 29 lines). No mathematics is written, repaired or re-derived: statement and proof are the article's own text line for line. Required context retained uncounted: sph (lines 241--245). Every definition, display and label of those lines is retained unchanged. Omitted from this package and not redistributed: 1--240, 246--384, 388--388, 418--676. Nothing inside a retained excerpt is omitted or altered: every retained body is delivered exactly as the article's lines are written, so no omission receipt of any kind accompanies this package. Citations: the retained excerpts carry no citation command at all, so no bibliography entry of the article is transcribed and no reference is redistributed. Reference adaptation: none was needed and none was made; every reference inside the delivered unit resolves inside it, so no unresolved reference or citation is produced. Printed slips kept literal: no printed word, symbol, hypothesis or number of the counted statement or of its proof was altered. Layout and class: the article's own class is \texttt{article} with packages including graphicx, amssymb, amsmath, amsthm, dsfont, enumitem, babel, tocbibind, nicematrix, biblatex; the wrapper is the lane's pinned \texttt{amsart} editorial wrapper at 10pt on A4, which restates the theorem-like environments the retained text needs and adds the article's own macro definitions that the retained text needs (none), with extra wrapper packages (bm, bbm, dsfont, xspace, cancel, mathtools). The delivered numbering is the unit's own. Assets: no retained span contains a figure environment, a graphics-inclusion command, a PDF-page inclusion, a table or a TikZ picture; no image file is copied into this package, the package declares no asset dependency and no image file is redistributed with it. Licence: version arXiv:2508.02367v2 of this article is distributed under the Creative Commons Attribution 4.0 International licence, as recorded in the arXiv licence metadata for this exact version and shown on the abstract page https://arxiv.org/abs/2508.02367v2. A full rights notice is delivered at licenses/arxiv-2508.02367-CC-BY-4.0.txt. Characters: every retained excerpt is pure ASCII, so the delivered engine, XeLaTeX with its default Latin Modern text font, reports no missing character.
\begin{thm} \label{sph}
For any irreducible spherical representation $(\pi,H)$ of $G$ that is not $1$-dimensional, there exists $\alpha\in\mathbb{K}$, $\alpha\neq\pm 1$, such that the representation is equivalent to the left action of $G$ on the subspace of $C$ defined by $\{h\in C: Lh=\alpha\cdot h\}$. Conversely, for every $\alpha\in\mathbb{K}$, $\alpha\neq\pm 1$, there is an irreducible spherical representation of $G$ that results from the construction.

(Note: It is not guaranteed that all these representations preserve a finite-index sesquilinear form; in fact, we will prove later that this is only possible if $\alpha\in\mathbb{R}$)
\end{thm}

\begin{thm} \label{sph2}
If $\alpha\neq 1$ and $\alpha\neq -\frac{1}{s-1}$, then the natural analog of Theorem \ref{sph} holds: the representation $(\pi,H)$ is equivalent to the left action of $G$ on the space $\{h\in C: L_2 h=\alpha\cdot h\}$, which is spanned by the $G$-translates of $f$, and conversely for every $\alpha\in\mathbb{K}$, $\alpha\neq 1$, $\alpha\neq -\frac{1}{s-1}$, the construction gives an irreducible spherical representation of $G$.
\end{thm}

\begin{proof}
We can repeat the proof of Theorem \ref{sph}. The only part that needs to be changed (and that will lead to the exception for $\alpha=-\frac{1}{s-1}$) is when we show that the space $\{h\in C: L_2 h=\alpha\cdot h\}$ is generated by the $G$-translates of the radial function $f$. To show this in our case, we fix a function $h$ in the space and try to build it with scalar multiples of translates of $f$. \\

We start with $F=f\cdot h(o)$, which agrees with $h$ at $o$. Now we wish to make it agree with $h$ on $S_2$ as well. For every $w\in S_2$, let $g_w\in G$ send $o$ to $w$. We add to $F$ the function $$\sum_{v\in N(o)}{\sum_{w\in N(v)}{\pi(g_w)f\cdot(f(0)-f(2))^{-1}(h(w)-F(w))}}.$$ The sum of the coefficients is zero, which guarantees that $F$ is unchanged at $o$. For every $v\in N(o)$, let $R(v):=\sum_{w\in N(v)}{(h(w)-F(w))}$. The condition $L_2 h=\alpha\cdot h$, together with $L_2 F=\alpha\cdot F$, implies that $\sum_{v\in N(o)}{R(v)}=0$. \\

For every $v\in N(o)$ and for every $w\in N(v)$, the function at $w$ has been changed by $$f(0)(f(0)-f(2))^{-1}(h(w)-F(w))$$ due to the contribution from $\pi(g_w)f$, by $$\sum_{z\in N(v)\setminus\{w\}}{f(2)(f(0)-f(2))^{-1}(h(z)-F(z))}=f(2)(f(0)-f(2))^{-1}(R(v)-(h(w)-F(w)))$$ due to the contribution from $\pi(g_z)f$ with $z\in N(v)\setminus\{w\}$, and by $$\sum_{x\in N(o)\setminus\{v\}}{\sum_{y\in N(x)}{f(4)(f(0)-f(2))^{-1}(h(y)-F(y))}}=-f(4)(f(0)-f(2))^{-1}R(v)$$ due to the contribution from $\pi(g_y)f$ with $y\in N(x)$ where $x\in N(o)\setminus\{v\}$. \\

All this adds up to $h(w)-F(w)+(f(2)-f(4))(f(0)-f(2))^{-1}R(v)$. \\

To remove the extra term $T(v):=(f(2)-f(4))(f(0)-f(2))^{-1}R(v)$, we subtract from $F$ the function $$\sum_{v\in N(o)}{\sum_{w\in N(v)}{\pi(g_w)f\cdot(f(0)+(s-2)f(2)-(s-1)f(4))^{-1}T(v)}}.$$

We need to assume that $f(0)+(s-2)f(2)-(s-1)f(4)\neq 0$. At every $w$ the value will be changed by $$-f(0)(f(0)+(s-2)f(2)-(s-1)f(4))^{-1}T(v)$$ due to the contribution from $\pi(g_w)f$, by $$-(s-2)f(2)(f(0)+(s-2)f(2)-(s-1)f(4))^{-1}T(v)$$ due to the contribution from $\pi(g_z)f$ with $z\in N(v)\setminus\{w\}$, and by
\begin{align*}
&\sum_{x\in N(o)\setminus\{v\}}{-(s-1)f(4)(f(0)+(s-2)f(2)-(s-1)f(4))^{-1}T(x)}= \\[0.5em]
&=(s-1)f(4)(f(0)+(s-2)f(2)-(s-1)f(4))^{-1}T(v)
\end{align*} due to the contribution from $\pi(g_y)f$ with $y\in N(x)$ where $x\in N(o)\setminus\{v\}$. \\

All this adds up to $-T(v)$, so now the function $F$ agrees with $h$ on $S_2$ as well as at $o$. \\

Like in the proof of Theorem \ref{sph}, we can follow the same procedure to make the two functions also agree on $S_4$ (applying the procedure to all vertices in $S_2$ in place of $o$), then on $S_6$ and so on, getting them to agree everywhere in a finite number of steps since $h$ is invariant under an open compact subgroup of $G$. \\

Therefore, the theorem is proved in the non-transitive case whenever $f(0)+(s-2)f(2)-(s-1)f(4)\neq 0$. As $f(2)=\alpha$ and
\begin{align*}
f(4)&=\frac{r}
{r-1}\left(\left(\alpha-\frac{s-2}{r(s-1)}\right)f(2)-\frac{1}{r(s-1)}f(0)\right)= \\
&=\frac{r(s-1)\alpha^2-(s-2)\alpha-1}{(r-1)(s-1)},
\end{align*}
the condition is true if and only if $1+(s-2)\alpha-\frac{r(s-1)\alpha^2-(s-2)\alpha-1}{(r-1)}\neq 0$, which is equivalent to $(s-1)\alpha^2-(s-2)\alpha-1\neq 0$, i.~e.~$((s-1)\alpha+1)(\alpha-1)\neq 0$, which holds for all $\alpha\neq 1,-\frac{1}{s-1}$. \\
\end{proof}

\end{document}
